% =====================================================================
%  Betriebsanweisung (Gefahrstoff) Phrozen Water-Washable Rapid Black
%  Eigenes PDF: Betriebsanweisung_Resin.tex, Seite in der Einweisung über \BAseite
%  Grundlage: Kennzeichnung von Phrozen Water-Washable Resin Schwarz
%  (Rapid Black) laut EU-Händlerangabe (3djake.de, Stand 28.09.2026):
%  GHS05, GHS07, GHS08, GHS09, Signalwort Gefahr,
%  H315, H317, H318, H335, H360FD, H373, H410.
%  Das Sicherheitsdatenblatt des Herstellers muss im FabLab vorliegen.
% =====================================================================
\BAsetup{
  typ=gefahrstoff,
  titel={Phrozen Water-Washable Rapid Black -- flüssiges Photopolymer-Harz auf Acrylatbasis},
  kurztitel={Gefahrstoff: Phrozen Water-Washable Rapid Black},
  taetigkeit={Drucken, Waschen im Form Wash (Wasser), Nachhärten, Harz einfüllen, Entsorgung},
  nummer=BA-GS-01,
  stand=28.09.2026,
  verantwortlich={Name der verantwortlichen Person},
  signalwort=GEFAHR,
}

\BAkopf

% ---------- 1 Gefahren -----------------------------------------------
\BAblock{GEFAHREN FÜR MENSCH UND UMWELT}{GHS08,GHS05,GHS07,GHS09}{%
  \begin{itemize}
    \item \textbf{Kann die Fruchtbarkeit beeinträchtigen. Kann das Kind im
          Mutterleib schädigen} (H360FD).
    \item \textbf{Verursacht schwere Augenschäden} (H318).
    \item \textbf{Kann allergische Hautreaktionen verursachen} (H317). Eine
          Allergie gegen Acrylate bleibt ein Leben lang bestehen.
    \item Verursacht Hautreizungen (H315). Kann die Atemwege reizen (H335).
          Kann die Organe schädigen bei längerer oder wiederholter
          Exposition (H373).
    \item \textbf{Sehr giftig für Wasserorganismen} mit langfristiger
          Wirkung (H410).
    \item Das gilt auch für \textbf{Waschwasser} und für nicht nachgehärtete
          Druckteile.
  \end{itemize}}

% ---------- 2 Schutzmaßnahmen ----------------------------------------
\BAblock{SCHUTZMASSNAHMEN UND VERHALTENSREGELN}{M004,M009,P022}{%
  \begin{itemize}
    \item \textbf{Schwangere und Stillende dürfen nicht mit dem Harz
          arbeiten} (Mutterschutzgesetz).
    \item \textbf{Schutzbrille} und \textbf{Nitrilhandschuhe} tragen;
          verschmutzte Handschuhe sofort wechseln.
    \item Haut- und Augenkontakt vermeiden, Dämpfe nicht einatmen.
          Raum gut lüften, Drucker- und Flaschendeckel geschlossen halten.
    \item Teile erst nach Waschen und Nachhärten ohne Handschuhe anfassen.
    \item Nach der Arbeit Hände waschen. Essen und Trinken verboten.
    \item Lagerung unter Verschluss, dunkel bei 15--35\,\textcelsius{},
          Flasche dicht verschlossen, nie in Lebensmittelgefäße umfüllen.
  \end{itemize}}

% ---------- 3 Gefahrfall ---------------------------------------------
\BAblock{VERHALTEN IM GEFAHRFALL}{W001,F001}{%
  \begin{itemize}
    \item Verschüttetes Harz mit Papiertüchern oder Bindemittel aufnehmen,
          nicht in den Abfluss gelangen lassen, Betreuer informieren.
    \item Verschmutzte Kleidung sofort ausziehen.
    \item Brand: Löschen mit Wassersprühstrahl, CO\textsubscript{2} oder
          Pulver. Brandgase (CO, NO\textsubscript{x}) nicht einatmen.
          \textbf{Feuerwehr: 112}
  \end{itemize}}

% ---------- 4 Erste Hilfe --------------------------------------------
\BAblock{ERSTE HILFE}{E003,E011}{%
  \begin{itemize}
    \item \textbf{Augen:} sofort mindestens 10 Min.\ bei geöffnetem Lid mit
          Wasser spülen, Kontaktlinsen entfernen, \textbf{sofort Augenarzt}.
    \item \textbf{Haut:} mit viel Wasser und Seife waschen. Bei Rötung
          oder Ausschlag Arzt aufsuchen.
    \item \textbf{Einatmen:} an die frische Luft, bei Beschwerden Arzt.
    \item \textbf{Verschlucken:} Mund ausspülen, kein Erbrechen
          herbeiführen, Arzt.
    \item Dem Arzt das Sicherheitsdatenblatt oder die Flasche zeigen.
  \end{itemize}\vspace{1.5mm}
  \textbf{Notruf: 112}\qquad Giftnotruf München: 089\,/\,19240\qquad
  FabLab: 09131\,/\,85\,28013}

% ---------- 5 Entsorgung (mit Fußzeile) --------------------------------
\BAletzterblock{SACHGERECHTE ENTSORGUNG}{}{%
  \begin{itemize}
    \item Flüssiges Harz und \textbf{Waschwasser nie in den Abfluss}.
          Das Waschwasser bleibt im Form Wash.
    \item Harzreste, Tücher und Handschuhe mit UV-Licht vollständig
          aushärten, danach im Restmüll entsorgen.
    \item Altes Waschwasser entsorgen nur Betreuer: aushärten lassen,
          feste Reste abfiltern (Restmüll).
  \end{itemize}}
